\documentclass[11pt,a4paper]{article}
\usepackage[utf8]{inputenc}
\usepackage[T1]{fontenc}
\usepackage[margin=2.5cm]{geometry}
\usepackage{amsmath,amssymb,amsthm}
\usepackage{booktabs}
\usepackage{longtable}
\usepackage{hyperref}

\newtheorem{theorem}{Theorem}
\newtheorem{lemma}[theorem]{Lemma}
\newtheorem{corollary}[theorem]{Corollary}
\theoremstyle{definition}
\newtheorem{definition}[theorem]{Definition}
\theoremstyle{remark}
\newtheorem{remark}[theorem]{Remark}

\newcommand{\GF}{\mathrm{GF}}
\newcommand{\PG}{\mathrm{PG}}
\newcommand{\PGL}{\mathrm{PGL}}
\newcommand{\cro}{\mathrm{cr}}
\newcommand{\QMAX}{1050}

\title{NSUCRYPTO 2026, Round 2\\Problem 3: The BSW Conjecture}
\date{}

\begin{document}
\maketitle

\section*{Summary}

We prove Conjecture~1 for finitely many values of $q$ (the third option in the task), and add two structural
results toward the general case: a stability theorem, and an exact reduction of the open case to a single
``sign'' property of triples.

\begin{theorem}[computer-assisted]\label{thm:main}
Conjecture~1 holds for every prime power $q\equiv 3\pmod 4$ with $31<q\le \QMAX$.
\end{theorem}

Blokhuis, Seress and Wilbrink report that A.~Brouwer checked $43\le q\le131$ by computer \cite[\S3]{BSW}.
Theorem~\ref{thm:main} extends this to $q\le\QMAX$, which includes the prime powers $243=3^5$ and $343=7^3$. The
computation reproduces the known small cases: for each $q\in\{7,11,19,23,27,31\}$ there are partitions that are
not of the form $n/x$, so the hypothesis $q>31$ cannot be weakened. This is the list in \cite[\S3]{BSW}, and for
prime $q$ the number of configurations up to $\PGL(2,q)$ (namely $1,2,1,2,1$ for $q=7,11,19,23,31$) also agrees
with theirs. Table~\ref{tab:small} gives the exact counts for $q\le31$.

\begin{theorem}[stability; all odd $q\ge5$]\label{thm:stab}
Let $\psi$ satisfy the hypotheses of Conjecture~1 for some odd prime power $q\ge 5$, and let $n$ be a nonsquare.
If $x\psi(x)=n$ holds for more than $(q-7+2\sqrt q)/2$ elements $x\in\GF(q)^*$, then $x\psi(x)=n$ for all
$x\in\GF(q)^*$. More generally, see Theorem~\ref{thm:stab2}: a ``BSW partition'' that shares more than
$(q-3+2\sqrt q)/4$ pairs with the orbits of an elliptic involution is equal to that orbit partition.
\end{theorem}

The threshold of Theorem~\ref{thm:stab2} is attained: for $q=7,11,19$ there are non-standard BSW partitions sharing
$\lfloor(q-3+2\sqrt q)/4\rfloor$ pairs with an elliptic involution (Section~\ref{sec:stats}).

Theorem~\ref{thm:stab} reduces Conjecture~1 to showing that some $n$ satisfies $x\psi(x)=n$ for roughly
half of all $x\in\GF(q)^*$.

\begin{theorem}[reduction; all $q\equiv3\pmod4$]\label{thm:red}
Let $q\equiv 3\pmod 4$ and let $\psi$ satisfy the hypotheses of Conjecture~1. Then $\psi(x)=n/x$ for some
nonsquare $n$ if and only if
\begin{equation}\label{eq:E}
\chi\bigl((s-t)(\psi(s)-\psi(t))\bigr)=1\qquad\text{for all nonzero squares } s\neq t .
\end{equation}
More symmetrically (Theorem~\ref{thm:red2}): every compatible triple of pairs has a sign $\pm1$ that is symmetric
and $\PGL(2,q)$-invariant, and a BSW partition is standard if and only if all of its triples are positive, if
and only if all triples through one of its pairs are positive. In particular, a BSW partition with fewer than
$(q+1)/6$ negative triples is standard.
\end{theorem}

The ``if'' direction adapts the first half of the proof of \cite{BSW} for $q\equiv1\pmod 4$. Their second half
proves the analogue of \eqref{eq:E} by a parity count in the Paley graph. For $q\equiv3\pmod4$ the same count
in the Paley tournament is valid but gives no information (Section~\ref{sec:failed}). Condition \eqref{eq:E} is
therefore exactly what remains open. We could not prove it; Section~\ref{sec:failed} lists the approaches that
failed.

\section{Notation and reformulation}

Throughout, $q$ is an odd prime power, $\chi$ is the quadratic character of $\GF(q)$, and
$\PG(1,q)=\GF(q)\cup\{\infty\}$. A point is written in homogeneous coordinates $x=(x_1:x_2)$, with
$a\in\GF(q)$ identified with $(a:1)$ and $\infty=(1:0)$. Put $[x,y]=x_1y_2-x_2y_1$. Then $[a,b]=a-b$ for finite
$a,b$, and $[\infty,b]=1$, $[a,\infty]=-1$.

\begin{definition}
For two disjoint $2$-subsets $P=\{a,a'\}$ and $Q=\{b,b'\}$ of $\PG(1,q)$ put
\[
  \cro(P,Q)=\frac{[a,b]\,[a',b']}{[a,b']\,[a',b]}\in\GF(q)^* .
\]
Each of $a,a',b,b'$ occurs once in the numerator and once in the denominator, so $\cro$ does not depend on the
choice of homogeneous coordinates. Swapping $a\leftrightarrow a'$ or $b\leftrightarrow b'$ replaces $\cro$ by
$1/\cro$, and swapping $P\leftrightarrow Q$ leaves it unchanged (each bracket changes sign, and there are two
in the numerator and two in the denominator). Hence $\chi(\cro(P,Q))$ and the \emph{class}
$\{\cro,\cro^{-1}\}$ of the unordered configuration $\{P,Q\}$ are well defined. We call $P,Q$
\emph{compatible} if $\chi(\cro(P,Q))=-1$.

A \emph{BSW partition} is a partition $\Pi$ of $\PG(1,q)$ into $(q+1)/2$ pairs, any two of which are
compatible. It is \emph{standard} if it is the set of orbits of a fixed-point-free involution in $\PGL(2,q)$.
\end{definition}

Since $[a,b][a',b'][a,b'][a',b]=\cro(P,Q)\cdot([a,b'][a',b])^2$, we have
\begin{equation}\label{eq:prod}
  \chi\bigl([a,b][a,b'][a',b][a',b']\bigr)=\chi(\cro(P,Q)).
\end{equation}

\begin{lemma}\label{lem:dict}
Let $\psi:\GF(q)^*\to\GF(q)^*$ be a fixed-point-free involution, and let
$\Pi_\psi=\{\{0,\infty\}\}\cup\{\{x,\psi(x)\}:x\in\GF(q)^*\}$. Then
\begin{enumerate}
\item[(i)] $\psi$ satisfies the hypotheses of Conjecture~1 if and only if $\Pi_\psi$ is a BSW partition;
\item[(ii)] $\psi(x)=n/x$ for a nonsquare $n$ if and only if $\Pi_\psi$ is standard.
\end{enumerate}
Every BSW partition containing the pair $\{0,\infty\}$ is $\Pi_\psi$ for exactly one such $\psi$.
\end{lemma}

\begin{proof}
(i) We have $\cro(\{0,\infty\},\{x,x'\})=\dfrac{[0,x][\infty,x']}{[0,x'][\infty,x]}=\dfrac{x}{x'}$, and
$\chi(x/x')=\chi(xx')$. So compatibility of $\{0,\infty\}$ with every other pair is exactly the requirement
$\chi(\psi(x))=-\chi(x)$. For $x,y\in\GF(q)^*$ with $y\notin\{x,\psi x\}$, the pairs $\{x,\psi x\}$ and
$\{y,\psi y\}$ are disjoint, and by \eqref{eq:prod} their compatibility is exactly the displayed condition of
Conjecture~1. The last sentence is clear: the pairs other than $\{0,\infty\}$ partition $\GF(q)^*$.

(ii) If $\psi(x)=n/x$ with $n$ a nonsquare, then $\sigma(x)=n/x$ (with $0\leftrightarrow\infty$) lies in
$\PGL(2,q)$. It is an involution, and it has no fixed point since $x^2=n$ has no solution. Its orbits form
$\Pi_\psi$. Conversely, let $\Pi_\psi$ be the orbit set of a fixed-point-free involution $\sigma(x)=(ax+b)/(cx+d)$.
Then $\sigma$ swaps $0$ and $\infty$. From $\sigma(0)=\infty$ we get $d=0$, and from $\sigma(\infty)=0$ we get
$a=0$. So $\sigma(x)=n/x$ with $n=b/c$, and $n$ is a nonsquare because $\sigma$ has no fixed point.
\end{proof}

\begin{lemma}\label{lem:stdbsw}
For every nonsquare $n$, the orbit set of $\sigma(x)=n/x$ is a BSW partition. Hence every standard partition
is a BSW partition.
\end{lemma}

\begin{proof}
The pair $\{0,\infty\}$ is compatible with $\{x,n/x\}$ because $\chi(x\cdot n/x)=\chi(n)=-1$. For two distinct finite
orbits $\{x,n/x\}$ and $\{y,n/y\}$, direct expansion gives
\[
(x-y)\Bigl(x-\frac ny\Bigr)\Bigl(\frac nx-y\Bigr)\Bigl(\frac nx-\frac ny\Bigr)
=n\left(\frac{(x-y)(xy-n)}{xy}\right)^{2},
\]
which is nonzero because the orbits are disjoint. Its character is $\chi(n)=-1$. By \eqref{eq:prod} the two
orbits are compatible. A standard partition is the orbit set of $g\sigma g^{-1}$ for some $g\in\PGL(2,q)$ and
some $\sigma$ as above: move one orbit to $\{0,\infty\}$ and apply Lemma~\ref{lem:dict}(ii). So it is BSW by
Lemma~\ref{lem:pgl}(i) below.
\end{proof}

\begin{lemma}\label{lem:pgl}
Let $g\in\PGL(2,q)$.
\begin{enumerate}
\item[(i)] $\cro(gP,gQ)=\cro(P,Q)$. Hence $g\Pi$ is a BSW partition if and only if $\Pi$ is, and $g$ maps
standard partitions to standard partitions.
\item[(ii)] If two unordered configurations $\{P,Q\}$ and $\{P',Q'\}$ of disjoint pairs have the same class,
then some $h\in\PGL(2,q)$ maps $\{P,Q\}$ to $\{P',Q'\}$.
\item[(iii)] For disjoint pairs $P,Q$, at most one involution of $\PGL(2,q)$ has both $P$ and $Q$ among its orbits.
\end{enumerate}
\end{lemma}

\begin{proof}
(i) For $g$ represented by a matrix $M$, $[Mx,My]=\det(M)[x,y]$. The factors $\det M$ cancel in $\cro$. If
$\Pi$ is the orbit set of $\sigma$, then $g\Pi$ is the orbit set of $g\sigma g^{-1}$, which is again a
fixed-point-free involution.

(ii) Write $P=\{a,a'\}$, $Q=\{b,b'\}$, $P'=\{\alpha,\alpha'\}$, $Q'=\{\beta,\beta'\}$. Relabelling within $P'$
if necessary, we may assume $\cro(P,Q)=\cro(P',Q')$, computed with these labels. $\PGL(2,q)$ is sharply
$3$-transitive on $\PG(1,q)$, so there is a unique $h$ with $h(a)=\alpha$, $h(a')=\alpha'$, $h(b)=\beta$. By
(i), $\cro(\{\alpha,\alpha'\},\{\beta,h(b')\})=\cro(\{\alpha,\alpha'\},\{\beta,\beta'\})$. For fixed distinct
$\alpha,\alpha',\beta$, the map $y\mapsto [\alpha,\beta][\alpha',y]/([\alpha,y][\alpha',\beta])$ is injective
on $\PG(1,q)\setminus\{\alpha,\alpha',\beta\}$: it is a M\"obius transformation in $y$ with nonzero determinant
$[\alpha,\beta][\alpha',\beta]\cdot[\alpha,\alpha']$ up to a nonzero scalar factor. Hence $h(b')=\beta'$.

(iii) If $\sigma,\sigma'$ are two such involutions, then $\sigma\sigma'$ fixes $a$, $a'$ and $b$, so it is the
identity and $\sigma=\sigma'^{-1}=\sigma'$.
\end{proof}

\begin{lemma}\label{lem:classes}
Let $q\equiv3\pmod4$ and let $g$ be a primitive element of $\GF(q)$. The class of every compatible
configuration is $\{g^L,g^{-L}\}$ for a unique odd \emph{key} $L\in\{1,3,\dots,(q-1)/2\}$. The class of $\{\{0,\infty\},\{1,g^L\}\}$ has key $L$.
In a standard partition, the configurations formed with $\{0,\infty\}$ realise every key.
\end{lemma}

\begin{proof}
A nonsquare $c$ is $g^{L}$ with $L$ odd, $0<L<q-1$, and $c^{-1}=g^{q-1-L}$. Since $q-1\equiv2\pmod4$,
$q-1-L$ is also odd. If $L\ne(q-1)/2$, exactly one of $L$ and $q-1-L$ is less than $(q-1)/2$; if $L=(q-1)/2$, i.e.\
$c=-1$, both equal $(q-1)/2$. So $\min(L,q-1-L)$ is the unique key. By the computation in the proof of Lemma~\ref{lem:dict}(i), $\cro(\{0,\infty\},\{1,g^L\})=g^{-L}$. In the standard
partition of $x\mapsto n/x$ we get $\cro(\{0,\infty\},\{x,n/x\})=x^2/n$, which runs over all nonsquares as
$x$ runs over $\GF(q)^*$.
\end{proof}

\section{A stability theorem}

\begin{lemma}\label{lem:hasse}
Let $q$ be odd and $F\in\GF(q)[X]$ monic of degree $4$ with four distinct roots in $\GF(q)$. Then
$T=\sum_{x\in\GF(q)}\chi(F(x))$ satisfies $|T+1|\le2\sqrt q$.
\end{lemma}

\begin{proof}
Write $F(X)=(X-r)(X-r_2)(X-r_3)(X-r_4)$. The map $u\mapsto r+1/u$ is a bijection from $\GF(q)^*$ onto
$\GF(q)\setminus\{r\}$, and $F(r+1/u)=u^{-4}G(u)$ with $G(U)=\prod_{i=2}^4\bigl((r-r_i)U+1\bigr)$. As $u^{-4}$ is a
nonzero square, $F(r)=0$ and $G(0)=1$, we get $T=\sum_{u\ne0}\chi(G(u))=\sum_{u\in\GF(q)}\chi(G(u))-1$. The cubic $G$ has
leading coefficient $c=\prod_i(r-r_i)\ne0$ and distinct roots $-1/(r-r_i)$, so $H(V)=c^2G(V/c)$ is a monic cubic with
distinct roots and $\chi(H(cu))=\chi(G(u))$. Since $q$ is odd, $y^2=H(v)$ is an elliptic curve $E$, and
$\#E(\GF(q))=1+\sum_v(1+\chi(H(v)))=q+1+\sum_v\chi(H(v))$. Hasse's theorem \cite[Theorem~V.1.1]{Silverman}
gives $|\sum_v\chi(H(v))|\le2\sqrt q$, and $\sum_v\chi(H(v))=\sum_u\chi(G(u))=T+1$.
\end{proof}

\begin{theorem}\label{thm:stab2}
Let $q\ge5$ be an odd prime power, $\Pi$ a BSW partition of $\PG(1,q)$, and $\sigma\in\PGL(2,q)$ a
fixed-point-free involution. If more than $(q-3+2\sqrt q)/4$ pairs of $\Pi$ are orbits of $\sigma$, then $\Pi$ is
the orbit set of $\sigma$.
\end{theorem}

\begin{proof}
Suppose some $\ell\in\Pi$ is not a $\sigma$-orbit. The $\sigma$-orbits partition $\PG(1,q)$ into
$(q+1)/2\ge3$ pairs, and $\ell$ meets at most two of them, so some $\sigma$-orbit $O$ is disjoint from $\ell$.
Choose $h\in\PGL(2,q)$ with $h(O)=\{0,\infty\}$ (possible by $2$-transitivity). Replace $\Pi,\sigma,\ell$ by
$h\Pi$, $h\sigma h^{-1}$, $h\ell$. By Lemma~\ref{lem:pgl}(i) the hypotheses and the conclusion are unchanged.
Now $\sigma$ swaps $0$ and $\infty$, so, as in the proof of Lemma~\ref{lem:dict}(ii), $\sigma(x)=n/x$ with $n$ a
nonsquare. Also $\ell=\{b,b'\}$ with $b,b'\in\GF(q)^*$ distinct, and $bb'\ne n$ because $\ell$ is not a
$\sigma$-orbit.

Every $\sigma$-orbit lying in $\Pi$ is disjoint from $\ell$ and compatible with it. We count the
$\sigma$-orbits with these two properties; the orbit $\{0,\infty\}$ is treated below. For $x\in\GF(q)^*$ with
$\{x,n/x\}\cap\{b,b'\}=\emptyset$, equation \eqref{eq:prod} gives
\[
\chi\bigl(\cro(\ell,\{x,n/x\})\bigr)=\chi\bigl((x-b)(x-b')(\tfrac nx-b)(\tfrac nx-b')\bigr)
=\chi(bb')\,\chi(F(x)),
\]
where $F(X)=(X-b)(X-b')(X-n/b)(X-n/b')$. Here we used $(\tfrac nx-b)(\tfrac nx-b')=\tfrac{bb'}{x^2}
(x-\tfrac nb)(x-\tfrac n{b'})$. The four roots of $F$ are distinct: $b\ne b'$; $b=n/b$ would give $b^2=n$, which
is impossible; $b=n/b'$ would give $bb'=n$, which is excluded; the same holds with $b,b'$ exchanged; and
$n/b\ne n/b'$. All four roots are nonzero. The condition
$\{x,n/x\}\cap\{b,b'\}=\emptyset$ says exactly that $F(x)\neq0$.

Let $M$ be the number of $x\in\GF(q)^*$ with $F(x)\ne0$ and $\chi(F(x))=-\chi(bb')$. Each compatible orbit
$\{x,n/x\}$ with $x\neq0$ contributes exactly two such $x$, since $x\ne n/x$. Hence the number of compatible
finite orbits is $M/2$. Since $F(0)=n^2\ne0$ and $\chi(F(0))=1$,
\[
M=\sum_{\substack{x\in\GF(q)^*\\F(x)\neq0}}\frac{1-\chi(bb')\chi(F(x))}{2}
=\frac12\Bigl((q-5)-\chi(bb')\bigl(T-1\bigr)\Bigr),\qquad T=\sum_{x\in\GF(q)}\chi(F(x)).
\]
The orbit $\{0,\infty\}$ is compatible with $\ell$ iff $\chi(bb')=-1$, since $\cro(\{0,\infty\},\ell)=b/b'$. So the number
$N$ of $\sigma$-orbits disjoint from and compatible with $\ell$ is $N=M/2=(q-4-T)/4$ if $\chi(bb')=1$, and
$N=1+M/2=(q-2+T)/4$ if $\chi(bb')=-1$. By Lemma~\ref{lem:hasse}, $-1-2\sqrt q\le T\le2\sqrt q-1$, so in both
cases $N\le(q-3+2\sqrt q)/4$. At most $N$ orbits of $\sigma$ lie in $\Pi$, contrary to the hypothesis.

So every pair of $\Pi$ is a $\sigma$-orbit. Both partitions consist of $(q+1)/2$ pairs, so they are equal.
\end{proof}

\begin{proof}[Proof of Theorem~\ref{thm:stab}]
Let $\Pi=\Pi_\psi$ and $\sigma(x)=n/x$. A pair $\{x,\psi x\}$ is a $\sigma$-orbit iff $x\psi(x)=n$, and each
such pair accounts for two elements $x$. The pair $\{0,\infty\}$ is also a common orbit. If more than
$(q-7+2\sqrt q)/2$ elements satisfy $x\psi(x)=n$, then $\Pi$ and the orbit set of $\sigma$ share more than
$(q-3+2\sqrt q)/4$ pairs, and Theorem~\ref{thm:stab2} finishes the proof.
\end{proof}

\section{A reduction for $q\equiv3\pmod 4$}\label{sec:red}

Throughout this section $q\equiv3\pmod4$, so $\chi(-1)=-1$ and $q=p^m$ with $m$ odd (an even power of an odd prime is
$\equiv1\pmod 4$).

\begin{definition}
A \emph{compatible triple} is a set $\{R,P,Q\}$ of three pairwise disjoint, pairwise compatible pairs. Choose
$g\in\PGL(2,q)$ with $gR=\{0,\infty\}$. By Lemma~\ref{lem:pgl}(i), $gP$ and $gQ$ are compatible with $\{0,\infty\}$, so
by the computation $\cro(\{0,\infty\},\{x,x'\})=x/x'$ in the proof of Lemma~\ref{lem:dict}(i) each of them consists of a nonzero square
and a nonsquare (none of the points is $0$ or $\infty$, by disjointness). Write $gP=\{s,s'\}$ and $gQ=\{t,t'\}$ with
$s,t$ squares, and put
\[
\varepsilon(R;P,Q)=\chi\bigl((s-t)(s'-t')\bigr)\in\{\pm1\}.
\]
($s\neq t$ and $s'\ne t'$ because $gP\cap gQ=\emptyset$.)
\end{definition}

\begin{lemma}\label{lem:eps}
$\varepsilon(R;P,Q)$ does not depend on the choice of $g$, it is $\PGL(2,q)$-invariant, and it is symmetric in
$R,P,Q$. We write $\varepsilon(R,P,Q)$ and call the triple \emph{positive} or \emph{negative}.
\end{lemma}

\begin{proof}
\emph{Independence of $g$.} Any other choice is $hg$ with $h$ in the stabiliser of $\{0,\infty\}$, i.e.\
$h(x)=\lambda x$ or $h(x)=\lambda/x$ with $\lambda\ne0$. If $h(x)=\lambda x$ with $\lambda$ a square, the labels are
kept and the product is multiplied by $\lambda^2$. If $\lambda$ is a nonsquare, $h$ swaps squares and nonsquares, so
the new square points are $\lambda s',\lambda t'$; the product $(\lambda s'-\lambda t')(\lambda s-\lambda t)$ is again
$\lambda^2(s-t)(s'-t')$. If $h(x)=\lambda/x$ with $\lambda$ a square, the labels are kept (as $\chi(\lambda/x)=\chi(x)$)
and $(\tfrac\lambda s-\tfrac\lambda t)(\tfrac\lambda{s'}-\tfrac\lambda{t'})=\lambda^2\frac{(s-t)(s'-t')}{st\,s't'}$, where
$\chi(st)=\chi(s't')=1$. If $\lambda$ is a nonsquare, the labels swap and the same product appears. In all cases
$\chi$ of the product is unchanged. Invariance under $k\in\PGL(2,q)$ follows by using $gk^{-1}$ for $kR$.

\emph{Symmetry.} Swapping $P$ and $Q$ swaps $(s,s')$ with $(t,t')$, which does not change $(s-t)(s'-t')$. It remains
to show $\varepsilon(R;P,Q)=\varepsilon(P;R,Q)$; the two transpositions generate $S_3$. Take $R=\{0,\infty\}$,
$P=\{s,s'\}$, $Q=\{t,t'\}$ as above. The map $g(x)=(x-s)/(x-s')$ lies in $\PGL(2,q)$ and sends $P$ to $\{0,\infty\}$,
$\infty\mapsto1$ (a square) and $0\mapsto s/s'$ (a nonsquare). Hence $\varepsilon(P;R,Q)=\chi((1-u)(s/s'-u'))$, where
$u$ is the square and $u'$ the nonsquare among $g(t),g(t')$. Direct computation gives
\[
1-g(y)=\frac{s-s'}{y-s'},\qquad \frac{s}{s'}-g(y)=\frac{y(s-s')}{s'(y-s')}\qquad(y\in\{t,t'\}).
\]
Recall $\chi(g(y))=\chi((y-s)(y-s'))$, and by compatibility of $P,Q$ exactly one of $g(t),g(t')$ is a square.

\emph{Case A: $g(t)$ is a square.} Then
$\varepsilon(P;R,Q)=\chi\bigl(\frac{(s-s')^2\,t'}{s'(t-s')(t'-s')}\bigr)=\chi(t's')\,\chi\bigl((t-s')(t'-s')\bigr)
=\chi\bigl((t-s')(t'-s')\bigr)$, since $t',s'$ are nonsquares. Comparing with
$\varepsilon(R;P,Q)=\chi((s-t)(s'-t'))$, the ratio is
$\chi\bigl((t-s')(s-t)\bigr)\chi\bigl((t'-s')(s'-t')\bigr)=\chi\bigl((t-s')(s-t)\bigr)\cdot\chi(-1)$. In Case A,
$\chi((t-s)(t-s'))=1$, so $\chi((s-t)(t-s'))=\chi(-1)=-1$, and the ratio is $(-1)(-1)=1$.

\emph{Case B: $g(t')$ is a square.} Then $\varepsilon(P;R,Q)=\chi\bigl(\frac{(s-s')^2\,t}{s'(t'-s')(t-s')}\bigr)
=\chi(ts')\chi\bigl((t-s')(t'-s')\bigr)=-\chi\bigl((t-s')(t'-s')\bigr)$. The ratio to $\varepsilon(R;P,Q)$ is
$-\chi((t-s')(s-t))\chi(-1)=\chi((t-s')(s-t))$. In Case B, $\chi((t-s)(t-s'))=-1$ (the other point is the square),
so $\chi((s-t)(t-s'))=\chi(-1)\cdot(-1)=1$, and the ratio is $1$.
\end{proof}

We use the following consequence of theorems of Carlitz \cite{Carlitz} and McConnel \cite{McConnel} (also proved by
Bruen and Levinger \cite{BL}), in the form stated in \cite[p.~145, ``Result'']{BSW}:

\begin{theorem}[Carlitz--McConnel]\label{thm:CM}
Let $q=p^m$ be odd and $f:\GF(q)\to\GF(q)$ a function such that $(f(x)-f(y))/(x-y)$ is a nonsquare for all
$x\ne y$. Then $f(x)=a+bx^{p^j}$ with $0\le j<m$ and $b$ a nonsquare.
\end{theorem}

\begin{theorem}\label{thm:red2}
Let $q\equiv3\pmod4$ and let $\Pi$ be a BSW partition of $\PG(1,q)$. The following are equivalent:
(a) $\Pi$ is standard; (b) every triple of pairs of $\Pi$ is positive; (c) some $R\in\Pi$ lies only in positive
triples of $\Pi$. Moreover, for $\Pi=\Pi_\psi$ and $R=\{0,\infty\}$, condition (c) is exactly \eqref{eq:E}.
\end{theorem}

\begin{proof}
Any three pairs of a BSW partition form a compatible triple.

(a)$\Rightarrow$(b). Let $\Pi$ be the orbit set of a fixed-point-free involution $\sigma$, and $\{R,P,Q\}\subseteq\Pi$.
Take $g$ with $gR=\{0,\infty\}$. Then $g\sigma g^{-1}$ swaps $0,\infty$, so it is $x\mapsto n/x$ with $n$ a nonsquare
(proof of Lemma~\ref{lem:dict}(ii)), and $gP=\{s,n/s\}$, $gQ=\{t,n/t\}$ with $s,t$ squares. So
$\varepsilon=\chi\bigl((s-t)(\tfrac ns-\tfrac nt)\bigr)=\chi\bigl(-n(s-t)^2/(st)\bigr)=\chi(-1)\chi(n)=1$.

(b)$\Rightarrow$(c) is trivial.

(c)$\Rightarrow$(a). By Lemma~\ref{lem:pgl}(i) and Lemma~\ref{lem:eps} we may move $R$ to $\{0,\infty\}$, so
$\Pi=\Pi_\psi$. The pairs of $\Pi$ other than $\{0,\infty\}$ are $\{s,\psi(s)\}$ with $s$ a square, so positivity of all
triples through $\{0,\infty\}$ is exactly \eqref{eq:E}; this also proves the last sentence of the theorem. Define
$f(0)=0$ and $f(x)=1/\psi(x)$ for $x\ne0$. We show that $D(x,y)=(f(x)-f(y))/(x-y)$ is a nonsquare for all $x\ne y$; as
$D$ is symmetric we may order $x,y$. Recall $\chi(1/z)=\chi(z)$, and that $\chi(\psi(s)\psi(t))=\chi(st)=1$ for squares
$s,t$.
\begin{itemize}
\item $y=0$: $D=1/(x\psi(x))$, and $\chi(x\psi(x))=-1$.
\item $y=\psi(x)$: $D=\dfrac{1/\psi(x)-1/x}{x-\psi(x)}=\dfrac1{x\psi(x)}$, a nonsquare.
\item $x=s$, $y=t$ distinct squares: $D=\dfrac{\psi(t)-\psi(s)}{\psi(s)\psi(t)(s-t)}$, so
$\chi(D)=\chi(-1)\chi\bigl((s-t)(\psi(s)-\psi(t))\bigr)=-1$ by \eqref{eq:E}.
\item $x=\psi(s)$, $y=\psi(t)$ with $s\ne t$ squares: $f(x)=1/s$, $f(y)=1/t$, so
$D=\dfrac{t-s}{st(\psi(s)-\psi(t))}$ and $\chi(D)=\chi(-1)\chi\bigl((s-t)(\psi(s)-\psi(t))\bigr)=-1$.
\item $x=s$ a square, $y=\psi(t)$ with $t\ne s$ a square: $D=\dfrac{t-\psi(s)}{t\,\psi(s)\,(s-\psi(t))}$, so
$\chi(D)=\chi(t\psi(s))\,\chi\bigl((t-\psi(s))(s-\psi(t))\bigr)=-\chi\bigl((t-\psi(s))(s-\psi(t))\bigr)$. By the hypothesis
of Conjecture~1 the product $(s-t)(s-\psi(t))(\psi(s)-t)(\psi(s)-\psi(t))$ is a nonsquare, and by \eqref{eq:E}
$(s-t)(\psi(s)-\psi(t))$ is a square, so $(s-\psi(t))(\psi(s)-t)$ is a nonsquare and
$\chi\bigl((t-\psi(s))(s-\psi(t))\bigr)=\chi(-1)\cdot(-1)=1$. Hence $\chi(D)=-1$.
\end{itemize}
These cases cover all $x\ne y$, since every nonzero element is a square $s$ or a nonsquare $\psi(s)$. By
Theorem~\ref{thm:CM}, $f(x)=a+bx^{p^j}$ with $b$ a nonsquare and $0\le j<m$; $f(0)=0$ gives $a=0$, so
$\psi(x)=1/(bx^{p^j})$ for $x\ne0$. Then $x=\psi(\psi(x))=b^{p^j-1}x^{p^{2j}}$ for all $x\ne0$. Taking $x=1$ gives
$b^{p^j-1}=1$, hence $x^{p^{2j}}=x$ on $\GF(q)$, so $m\mid 2j$. As $m$ is odd, $m\mid j$, and $0\le j<m$ gives $j=0$.
So $\psi(x)=n/x$ with $n=1/b$ a nonsquare, and $\Pi$ is standard by Lemma~\ref{lem:dict}(ii).
\end{proof}

\begin{corollary}\label{cor:few}
If a BSW partition $\Pi$ ($q\equiv3\pmod4$) has fewer than $(q+1)/6$ negative triples, it is standard.
\end{corollary}

\begin{proof}
The negative triples contain at most $3\cdot\#\{\text{negative triples}\}<(q+1)/2=|\Pi|$ pairs, so some $R\in\Pi$ lies
in no negative triple. Apply Theorem~\ref{thm:red2}(c).
\end{proof}

\begin{proof}[Proof of Theorem~\ref{thm:red}]
By Lemma~\ref{lem:dict}, $\psi(x)=n/x$ iff $\Pi_\psi$ is standard, and by Theorem~\ref{thm:red2} (with
$R=\{0,\infty\}$) this holds iff \eqref{eq:E} holds. The remaining statements are Lemma~\ref{lem:eps},
Theorem~\ref{thm:red2} and Corollary~\ref{cor:few}.
\end{proof}

\begin{remark}
In the non-standard partitions for $q\le31$ the property fails massively: for every choice of $R$, between
$39\%$ and $100\%$ of the pairs $\{P,Q\}$ give negative triples (for $q=31$, $50$ or $66$ of the $105$). Standard
partitions have none. Numerically, $\varepsilon$ was also checked to be symmetric on $12{,}000$ random compatible
triples for $q\in\{19,43,103,227\}$, and the equivalence of \eqref{eq:E}, the nonsquare property of $f$, and
standardness was checked on all $244$ BSW partitions containing $\{0,\infty\}$ for the primes $q\equiv3\pmod4$
with $7\le q\le43$.
\end{remark}

\section{Computer-assisted proof of Theorem~\ref{thm:main}}

\subsection{Reduction to finitely many searches}

Fix $q\equiv3\pmod4$ and a primitive element $g$. For an odd key $K\le(q-1)/2$, let $\mathcal S_K$ be the set of
BSW partitions $\Pi$ that contain $\{0,\infty\}$ and $\{1,g^K\}$ and in which every two pairs have key
$\ge K$.

\begin{lemma}\label{lem:orderly}
\begin{enumerate}
\item[(i)] Every BSW partition of $\PG(1,q)$ is $\PGL(2,q)$-equivalent to a member of $\mathcal S_K$ for some
$K$.
\item[(ii)] If $q\equiv3\pmod 4$, then $\mathcal S_1$ contains exactly one standard partition, namely the orbits
of $x\mapsto g/x$, and $\mathcal S_K$ contains no standard partition for $K>1$.
\end{enumerate}
Consequently, if every $\mathcal S_K$ consists of standard partitions only, then Conjecture~1 holds for $q$.
\end{lemma}

\begin{proof}
(i) Let $K$ be the least key over all configurations of two pairs of $\Pi$, attained by $\{P,Q\}$. By
Lemma~\ref{lem:classes}, $\{\{0,\infty\},\{1,g^K\}\}$ also has key $K$. By Lemma~\ref{lem:pgl}(ii) some $h$ maps
$\{P,Q\}$ to it. By Lemma~\ref{lem:pgl}(i), $h\Pi$ is a BSW partition with the same multiset of keys, so
$h\Pi\in\mathcal S_K$.

(ii) Standard partitions are BSW partitions (Lemma~\ref{lem:stdbsw}). A standard partition realises key $1$ (Lemma~\ref{lem:classes}, applied after moving one of its pairs to
$\{0,\infty\}$; keys are invariant). So its least key is $1$, and it lies in no $\mathcal S_K$ with $K>1$. A
standard partition in $\mathcal S_1$ is the orbit set of an involution having $\{0,\infty\}$ and $\{1,g\}$ as
orbits. By Lemma~\ref{lem:pgl}(iii) this involution is unique, and $x\mapsto g/x$ is one.

For the final claim, let $\psi$ satisfy the hypotheses of Conjecture~1. Then $\Pi_\psi$ is a BSW partition
(Lemma~\ref{lem:dict}). By (i), $h\Pi_\psi\in\mathcal S_K$ for some $h$ and $K$. By assumption $h\Pi_\psi$ is
standard, so $\Pi_\psi$ is standard by Lemma~\ref{lem:pgl}(i), and $\psi(x)=n/x$ by Lemma~\ref{lem:dict}(ii).
\end{proof}

\subsection{The search}

For a member of $\mathcal S_K$, every pair other than $\{0,\infty\}$ and $\{1,g^K\}$ is $\{s,t\}$ with $s$ a
nonzero square, $s\ne1$, and $t$ a nonsquare, $t\ne g^K$ (Lemma~\ref{lem:dict}(i)). So a member of
$\mathcal S_K$ is a perfect matching between the $N=(q-3)/2$ ``rows'' $s$ and the $N$ ``columns'' $t$ such that:
\begin{itemize}
\item each chosen pair $\{s,t\}$ is \emph{allowed}: its configurations with $\{0,\infty\}$ and with
$\{1,g^K\}$ have key $\ge K$ (keys are only defined for compatible configurations, so this includes
compatibility);
\item every two chosen pairs form a configuration of key $\ge K$.
\end{itemize}
We enumerated all such matchings exhaustively by backtracking, for every prime power $q\equiv3\pmod4$ with
$43\le q\le1050$ and every odd key $K\le(q-1)/2$. Each member of $\mathcal S_K$ is reported exactly once. By
Lemma~\ref{lem:orderly}(i) this covers every BSW partition up to $\PGL(2,q)$, so the results below prove
Theorem~\ref{thm:main}.

\subsection{Verification}

The following independent checks were performed.
\begin{enumerate}
\item \textbf{Positive control.} By Lemma~\ref{lem:orderly}(ii), the search must report the standard partition
$x\mapsto g/x$ in run $K=1$. It did so for every $q$ in the range.
\item \textbf{Output check.} Every reported partition was re-checked by an independent program. It rebuilds
$\psi$ from the printed pairs, checks that $\psi$ is an involution of $\GF(q)^*$ with $\chi(\psi(x))=-\chi(x)$,
checks the four-factor condition of Conjecture~1 for all $x,y$ directly, and decides whether $x\psi(x)$ is
constant. For prime powers it also re-verifies that the printed polynomial defines a field.
\item \textbf{Independent re-run by SAT.} For every prime $q\equiv3\pmod4$ with $43\le q\le131$, all BSW
partitions containing $\{0,\infty\}$ were also enumerated with a SAT encoding that uses no keys and no other
symmetry reduction: one Boolean variable per pair $\{s,t\}$, exactly-one constraints on rows and columns, a
binary clause for every incompatible configuration, and blocking clauses to enumerate all models. This finds
exactly $(q-1)/2$ partitions, all standard. Together with Lemma~\ref{lem:dict} this proves the conjecture for
these $q$ without using the search program.
\item \textbf{Small $q$.} For $q\le31$, the set of partitions containing $\{0,\infty\}$ found by the search was
reproduced by two further independent enumerations (including $q=27$). The $q\le31$ with non-standard partitions
are exactly $7,11,19,23,27,31$ (Table~\ref{tab:small}).
\item \textbf{Third method.} An independently written exhaustive search, with its own field construction, using
no logarithms, no keys and no orderly reduction (it fixes only $\{0,\infty\}$ and branches over the partner of
$1$), reproduces the main search for all $29$ prime powers $q\equiv3\pmod4$ with $43\le q\le311$ (including
$243$): exactly $(q-1)/2$ BSW partitions containing $\{0,\infty\}$, all standard. This re-proves
Theorem~\ref{thm:main} for these $q$ independently of the main search, and for $q\le31$ its counts equal
Table~\ref{tab:small}.
\end{enumerate}
For $311<q\le700$ the result rests on the main search, supported by the positive control and the output check.

\begin{table}[ht]
\centering
\caption{All BSW partitions of $\PG(1,q)$ for $7\le q\le 31$, $q\equiv 3\pmod 4$ ($q=3$ is trivial: the two pairs always form the standard partition). ``Containing $\{0,\infty\}$'' is the
count found by the independent exhaustive enumerations (SAT for prime $q$, brute force otherwise); ``total'' is $q$ times that count, since each partition has $(q+1)/2$ of the $q(q+1)/2$ pairs and,
by transitivity, every pair lies in equally many partitions.}
\label{tab:small}
\begin{tabular}{rrrrr}
\toprule
$q$ & containing $\{0,\infty\}$ & of which standard & total & non-standard total\\
\midrule
7 & 5 & 3 & 35 & 14\\
11 & 12 & 5 & 132 & 77\\
19 & 54 & 9 & 1026 & 855\\
23 & 121 & 11 & 2783 & 2530\\
27 & 104 & 13 & 2808 & 2457\\
31 & 31 & 15 & 961 & 496\\
\bottomrule
\end{tabular}
\end{table}

\subsection{Results}

For every prime power $q\equiv3\pmod4$ with $43\le q\le1050$ (Table~\ref{tab:big}), every run $K$ completed.
Run $K=1$ reported exactly one partition, the standard one, and every other run reported none. The output check
confirmed every output. By Lemma~\ref{lem:orderly} this proves Theorem~\ref{thm:main}. The range
$700<q\le1050$ was computed by a distributed version of the same search (Section~\ref{sec:dist}). All source
code, logs and computed outputs are available upon request.

\begin{longtable}{rrr}
\caption{Searches for $43\le q\le1050$: number of runs $K$ (all of them, $(q+1)/4$) and partitions reported (all
runs together; in every case the single standard partition).}\label{tab:big}\\
\toprule
$q$ & runs & reported\\
\midrule
\endfirsthead
\toprule
$q$ & runs & reported\\
\midrule
\endhead
43 & 11 & 1 \\
47 & 12 & 1 \\
59 & 15 & 1 \\
67 & 17 & 1 \\
71 & 18 & 1 \\
79 & 20 & 1 \\
83 & 21 & 1 \\
103 & 26 & 1 \\
107 & 27 & 1 \\
127 & 32 & 1 \\
131 & 33 & 1 \\
139 & 35 & 1 \\
151 & 38 & 1 \\
163 & 41 & 1 \\
167 & 42 & 1 \\
179 & 45 & 1 \\
191 & 48 & 1 \\
199 & 50 & 1 \\
211 & 53 & 1 \\
223 & 56 & 1 \\
227 & 57 & 1 \\
239 & 60 & 1 \\
243 & 61 & 1 \\
251 & 63 & 1 \\
263 & 66 & 1 \\
271 & 68 & 1 \\
283 & 71 & 1 \\
307 & 77 & 1 \\
311 & 78 & 1 \\
331 & 83 & 1 \\
343 & 86 & 1 \\
347 & 87 & 1 \\
359 & 90 & 1 \\
367 & 92 & 1 \\
379 & 95 & 1 \\
383 & 96 & 1 \\
419 & 105 & 1 \\
431 & 108 & 1 \\
439 & 110 & 1 \\
443 & 111 & 1 \\
463 & 116 & 1 \\
467 & 117 & 1 \\
479 & 120 & 1 \\
487 & 122 & 1 \\
491 & 123 & 1 \\
499 & 125 & 1 \\
503 & 126 & 1 \\
523 & 131 & 1 \\
547 & 137 & 1 \\
563 & 141 & 1 \\
571 & 143 & 1 \\
587 & 147 & 1 \\
599 & 150 & 1 \\
607 & 152 & 1 \\
619 & 155 & 1 \\
631 & 158 & 1 \\
643 & 161 & 1 \\
647 & 162 & 1 \\
659 & 165 & 1 \\
683 & 171 & 1 \\
691 & 173 & 1 \\
719 & 180 & 1 \\
727 & 182 & 1 \\
739 & 185 & 1 \\
743 & 186 & 1 \\
751 & 188 & 1 \\
787 & 197 & 1 \\
811 & 203 & 1 \\
823 & 206 & 1 \\
827 & 207 & 1 \\
839 & 210 & 1 \\
859 & 215 & 1 \\
863 & 216 & 1 \\
883 & 221 & 1 \\
887 & 222 & 1 \\
907 & 227 & 1 \\
911 & 228 & 1 \\
919 & 230 & 1 \\
947 & 237 & 1 \\
967 & 242 & 1 \\
971 & 243 & 1 \\
983 & 246 & 1 \\
991 & 248 & 1 \\
1019 & 255 & 1 \\
1031 & 258 & 1 \\
1039 & 260 & 1 \\
\bottomrule
\end{longtable}

\subsection{Distributed computation}\label{sec:dist}

The range $700<q\le1050$ consists of the $25$ primes $q\equiv3\pmod4$ from $719$ to $1039$; there is no proper
prime power $q\equiv3\pmod4$ in this range. A distributed version of the same search computed all runs $K$ in
this range, and the merged outputs passed the output check described above. All $25$ values of $q$ are complete,
and each has exactly one solution, the standard one (Table~\ref{tab:big}).

\section{Remarks}

\begin{itemize}
\item \emph{The sporadic case $q=31$.} The $496$ non-standard partitions for $q=31$ form one $\PGL(2,31)$-orbit
(stabiliser of order $60$). Each of them meets $15$ elliptic
involutions in $4$ pairs and $30$ in $2$ pairs. So every involution contains at most $4$ of its $16$ pairs, far
below the bound $(28+2\sqrt{31})/4\approx9.78$ of Theorem~\ref{thm:stab2}.
\end{itemize}

\section{Statistics of the compatibility relation}\label{sec:stats}

These experiments are heuristic; no theorem depends on them.
\begin{itemize}
\item \emph{Raw search tree.} Fix $\{0,\infty\}$, order the squares $s_1,\dots,s_m$ ($m=(q-1)/2$), and let $N_d$ be
the number of tuples of distinct nonsquares $(n_1,\dots,n_d)$ with $\{s_i,n_i\}$ pairwise compatible. Positive
control: $N_m$ equals the counts of Table~\ref{tab:small}
($5,12,54,121,31$ for $q=7,\dots,31$). If compatibilities were independent fair coins,
$\mathbb E N_d=m(m-1)\cdots(m-d+1)\,2^{-\binom d2}$. Over $3$ square orders per $q$, the ratio $(N_d-m)/\mathbb E N_d$
at $d=4$ is $1.21$ for $q=43$, $1.06$ for $q=107$, $1.02$ for $q=227$ and $1.017$ for $q=503$. Non-standard branches
die $0$ to $7$ levels below the depth where $\mathbb E N_d<1$. The deep survivors are near-standard: for $q=103$ and
$d\ge12$ almost all have every pair but one in a single involution, the situation of Theorem~\ref{thm:stab2}.
\item \emph{Local independence.} For $k=3,4$ random disjoint pairs ($4\cdot10^6$ samples per $q$), the
$2^{\binom k2}$ compatibility patterns are uniform within sampling error for $q\ge10007$ ($\chi^2$ $p$-values
$0.3$--$0.94$), and nearly so for $q=2003,5003$ (smallest $p$-value $0.014$). For smaller $q$ the deviations are of order $1/q$, not $1/\sqrt q$. For example, the probability
that two pairs are compatible is about $\frac12+\frac1{2q}$. The sign $\varepsilon$ of
Lemma~\ref{lem:eps} is positive on random compatible triples with frequency $0.496$--$0.503$ for
$q\in\{43,227,503,1031,5003,20011\}$, but $0.492$ for $q=103$ ($p=0.001$).
\item \emph{The exceptions are symmetric, not random.} The independent-coin expectation of the number of BSW
partitions containing $\{0,\infty\}$ is $m!\,2^{-\binom m2}$: $0.75$, $0.12$, $5\cdot10^{-6}$, $10^{-9}$ and
$3\cdot10^{-20}$ for $q=7,11,19,23,31$. The observed numbers of non-standard ones are $2,7,45,110,16$. Every
non-standard $\PGL(2,q)$-orbit has a nontrivial stabiliser, of order $24$; $60,24$; $8$; $8,12$; $60$.
\item \emph{Sharpness of Theorem~\ref{thm:stab2}.} The maximum number of shared pairs between a non-standard BSW
partition and an elliptic involution is $2,3,6,4,4$ for $q=7,11,19,23,31$. The thresholds are $2.32$, $3.66$, $6.18$,
$7.40$ and $9.78$, so the theorem is sharp for $q=7,11,19$. In the proof, the maximum of $N$ over $\ell$ equals
$\lfloor(q-3+2\sqrt q)/4\rfloor$ for every tested prime $q\in\{43,59,83,103,131,163,199,251,307\}$.
\end{itemize}

\section{Approaches that failed}\label{sec:failed}

\begin{itemize}
\item \emph{Forcing a large common involution.} Theorem~\ref{thm:stab2} would prove the conjecture for all
large $q$, given a lemma that some elliptic involution shares more than $(q-3+2\sqrt q)/4$ pairs with $\Pi$. Pure
counting gives only $\sum_z t_z=\sum_z t_z(t_z-1)=(q^2-1)/4$, where $t_z$ is the number of pairs of $\Pi$ in the
involution $z$. This is compatible with all $t_z\le2$, so it gives nothing.
\item \emph{Parity linearisation.} For $q\equiv3\pmod4$, two pairs $P,Q$ are compatible iff an odd number of the four
pairs $(x,y)\in P\times Q$ have $\chi([y,x])=1$, by \eqref{eq:prod}. This turns the condition into linear algebra over
$\GF(2)$, but we found no way to exploit it.
\item \emph{Odd character sums.} $\Pi$ is BSW iff each $g_P(x)=\chi([x,a][x,a'])$ (bracket convention, so
$\infty$ is allowed; $P=\{a,a'\}$) is odd under the pairing
outside $P$. Hence for every three pairs, $\sum_{x\in\PG(1,q)}g_P(x)g_Q(x)g_R(x)=0$ (the terms for $x$ and its partner cancel),
i.e.\ the genus-$2$ curve $y^2=f_Pf_Qf_R$ has Frobenius trace $0$ (checked on $10{,}112$ triples). We found no way to turn these identities into a contradiction.
\item \emph{The parity count of \cite{BSW}.} For $q\equiv1\pmod4$, BSW prove the analogue of \eqref{eq:E} by counting,
for the $5$-set $X=\{0,a,b,c,d\}$ formed by $0$ and two pairs, the vertices joined to exactly $i$ points of $X$ in the
(strongly regular) Paley graph. For $q\equiv3\pmod4$ we repeated the count in the Paley tournament ($z\to u$ iff
$\chi(u-z)=1$). For $v\notin X$ let $i(v)$ be the number of $u\in X$ with $\chi(u-v)=1$, and $x_i$ the number of $v$
with $i(v)=i$. Since any two vertices have exactly $(q-3)/4$ common in-neighbours (a direct character-sum
computation), one gets $\sum x_i=q-5$, $\sum ix_i=5\frac{q-1}2-10$ and $\sum\binom i2x_i=10\frac{q-3}4-S$, where $S$ is the
number of transitive triples of the sub-tournament on $X$. Compatibility gives, as in \cite{BSW}, $\sum_{i\text{ odd}}x_i=
\frac{q-5}2$. Hence $2x_3+4x_4+8x_5=\frac{3q-5}2-S$, so $S$ must be even. We verified this identity on actual partitions
($q=19,23,31,43$). But enumerating all $32$ sub-tournaments allowed by the
known arcs and the compatibility of the two pairs shows that $S$ is always even. So
the count gives no information.
\item \emph{Probabilistic and pseudorandom counting.} First- and second-moment counts with character sums bound
the size of families of pairwise compatible pairs (Cauchy--Schwarz with Lemma~\ref{lem:hasse} gives roughly
$k^2\lesssim q^2/2+2kq$). But standard partitions are such families of the largest possible size $(q+1)/2$, so no
bound valid for all compatible families can exclude non-standard ones. The first moments of
Section~\ref{sec:stats} are heuristics only.
\item \emph{Delsarte's linear programme.} $\PGL(2,q)$ acting on the $2$-subsets of $\PG(1,q)$ gives a commutative,
multiplicity-free association scheme with $(q+3)/2$ classes (identity, ``meet'', and the cross-ratio classes
$\{c,1/c\}$). A BSW partition is a clique avoiding the square classes and ``meet'', and it is orthogonal to the
point-incidence module. For $q=43,47$ the resulting LP leaves every nonsquare class free over $[0,\approx5.5]$, so it
does not force the inner distribution of the standard partition. Moreover,
standard partitions are Delsarte-slack, vanishing only on the point module, while the non-standard examples are
tight (for $q=31$ on $7$ eigenspaces). So Delsarte equality cannot single out the standard partitions.
\item \emph{A heuristic obstruction.} Non-standard partitions exist for $q=7,11,19,23,27,31$, which meet every
residue class $q\equiv3,7,11,15\pmod{16}$. So an argument that uses only identities valid for all $q$ in one of
these classes modulo $16$ (exact counts, parity) cannot prove Conjecture~1. This suggests, but does not prove, that a
proof must use the size of $q$, for instance through Weil-type bounds as in Theorem~\ref{thm:stab2}.
\end{itemize}

\begin{thebibliography}{9}
\bibitem{BSW} A.~Blokhuis, \'A.~Seress, H.~A.~Wilbrink, Characterization of complete exterior sets of conics,
\emph{Combinatorica} 12 (1992) 143--147.
\bibitem{Carlitz} L.~Carlitz, A theorem on permutations in a finite field, \emph{Proc. Amer. Math. Soc.} 11 (1960)
456--459.
\bibitem{McConnel} R.~McConnel, Pseudo-ordered polynomials over a finite field, \emph{Acta Arith.} 8 (1963) 127--151.
\bibitem{BL} A.~A.~Bruen, B.~Levinger, A theorem on permutations of a finite field, \emph{Canad. J. Math.} 25 (1973)
1060--1065.
\bibitem{Silverman} J.~H.~Silverman, \emph{The Arithmetic of Elliptic Curves}, 2nd ed., Springer, 2009.
\end{thebibliography}

\end{document}
